\subsection{Inference, Prediction, and Multiple Regression}
% Source integration: attached Chapter 14 screenshots, Sections 14.4.2--14.5.

Once the linear model is written in matrix form, standard errors for coefficients and predictions are consequences of linear transformations of random vectors.

\begin{theorem} \label{thm:linear-model-beta-hat-mean-covariance} \hlt{Mean and Covariance of Least-Squares Estimates}\\
Assume $Y=X\beta+e$, $\operatorname{rank}(X)=p$, $\mathbb{E}e=0$, and $\operatorname{Var}(e)=\sigma^2I_n$.  Then
\[
\mathbb{E}\widehat{\beta}=\beta,
\qquad
\operatorname{Var}(\widehat{\beta})=\sigma^2(X^tX)^{-1}.
\]
\begin{proof}
From Theorem \ref{thm:linear-model-normal-equations},
\[
\widehat{\beta}=(X^tX)^{-1}X^t(X\beta+e)
=\beta+(X^tX)^{-1}X^te.
\]
The expectation follows from $\mathbb{E}e=0$.  The covariance formula follows from Proposition \ref{prop:linear-transform-random-vector-moments}:
\[
\operatorname{Var}(\widehat{\beta})
=(X^tX)^{-1}X^t(\sigma^2I_n)X(X^tX)^{-1}
=\sigma^2(X^tX)^{-1}.
\]
\end{proof}
\end{theorem}

\begin{method} \label{meth:linear-model-coefficient-inference} \hlt{Coefficient Inference in a Normal Linear Model}\\
Let $C=(X^tX)^{-1}$ and estimate the standard error of $\widehat{\beta}_j$ by
\[
s_{\widehat{\beta}_j}=s\sqrt{c_{jj}}.
\]
If the errors are independent normal variables, then
\[
\frac{\widehat{\beta}_j-\beta_j}{s_{\widehat{\beta}_j}}\sim t_{n-p}.
\]
Thus a $100(1-\alpha)\%$ confidence interval is
\[
\widehat{\beta}_j\pm t_{n-p,1-\alpha/2}s_{\widehat{\beta}_j},
\]
and the test of $H_0:\beta_j=\beta_{j0}$ uses the statistic
\[
t=\frac{\widehat{\beta}_j-\beta_{j0}}{s_{\widehat{\beta}_j}}.
\]
Asymptotically, similar intervals are often used under weaker conditions by central-limit reasoning.
\end{method}

\begin{proposition} \label{prop:linear-model-mean-response-interval} \hlt{Mean Response at a New Design Point}\\
For a fixed predictor vector $x_0\in\R^p$, the estimated mean response is
\[
\widehat{\mu}_0=x_0^t\widehat{\beta}.
\]
Under the assumptions of Theorem \ref{thm:linear-model-beta-hat-mean-covariance},
\[
\operatorname{Var}(\widehat{\mu}_0)=\sigma^2x_0^t(X^tX)^{-1}x_0.
\]
Replacing $\sigma$ by $s$ gives the usual standard error $s_{\widehat{\mu}_0}$ and normal-model confidence interval
\[
\widehat{\mu}_0\pm t_{n-p,1-\alpha/2}s_{\widehat{\mu}_0}.
\]
\end{proposition}

\begin{corollary} \label{cor:linear-model-prediction-interval} \hlt{Prediction Interval for a New Observation}\\
If a new observation at $x_0$ is
\[
Y_0=x_0^t\beta+e_0,
\]
with $e_0$ independent of the original errors and variance $\sigma^2$, then
\[
\operatorname{Var}(Y_0-\widehat{\mu}_0)
=\sigma^2\left(1+x_0^t(X^tX)^{-1}x_0\right).
\]
Under normality, a $100(1-\alpha)\%$ prediction interval is
\[
\widehat{\mu}_0\pm t_{n-p,1-\alpha/2}s
\sqrt{1+x_0^t(X^tX)^{-1}x_0}.
\]
\end{corollary}

\begin{definition} \label{def:multiple-linear-regression-coefficients} \hlt{Multiple Regression Coefficients}\\
In the model
\[
Y_i=\beta_0+\beta_1x_{i1}+\cdots+\beta_{p-1}x_{i,p-1}+e_i,
\]
the coefficient $\beta_k$ is the change in the expected response associated with increasing $x_k$ by one unit while holding the other predictors fixed.  Polynomial regression is included by treating powers such as $x_i^2$ as additional columns of the design matrix.
\end{definition}

\begin{definition} \label{def:coefficient-of-determination-r-squared} \hlt{Coefficient of Determination}\\
For a model with an intercept, define
\[
\operatorname{TSS}=\sum_{i=1}^{n}(Y_i-\overline{Y})^2,
\qquad
R^2=1-\frac{\operatorname{RSS}}{\operatorname{TSS}}.
\]
Equivalently, $R^2$ is the squared sample correlation between $Y$ and $\widehat{Y}$.  It measures the fraction of response variability removed by the fitted model on the observed sample.
\end{definition}

\begin{remark} \label{rem:collinearity-and-coefficient-interpretation} \hlt{Collinearity and Interpretation}\\
When two predictors are strongly linearly related, the fitted mean may be stable while individual coefficient estimates have large standard errors.  In the catheter-length example from the source, height and weight predicted length similarly one at a time; fitting both together added little reduction in residual scale and made coefficient interpretation unstable.  Exact linear dependence would make $X$ rank deficient, so the normal equations would not have a unique coefficient vector.
\end{remark}

\begin{example} \label{ex:quadratic-term-as-regression-check} \hlt{Quadratic Terms as a Lack-of-Fit Check}\\
The source revisits a peak-area calibration fit by adding $x^2$ to the straight-line model.  The quadratic coefficient had a small $t$ statistic, so there was no evidence that the extra term was needed.  In contrast, for the stream-flow example, a quadratic term in depth was strongly significant and the residual plot no longer showed systematic curvature.  The lesson is diagnostic rather than automatic: an added term should be motivated by residual structure or subject-matter scale.
\end{example}
